\begin{abstract}
\textbf{Objective:}
Covariance-estimation fidelity and discrimination of unseen EEG participants
need not favour the same regularisation strength. We ask whether this
difference changes zero-target cross-subject decoding and conclusions about
spectral representations.
\textbf{Methods:}
Within a fixed covariance/tangent-space linear-discriminant family, we
compare source-participant macro-AUC selection, source covariance
predictive negative log-likelihood selection and analytic Ledoit--Wolf
shrinkage over the full scalar path. All reported target comparisons use
source-only selection. Experiments include prospectively held-out EEGMMIDB
participants, two separately pre-data-frozen BNCI datasets, an independently
frozen Shin2017A cohort and a distinct prespecified representation test.
Participant-level paired inference and method-specific multiplicity
correction are retained.
\textbf{Results:}
On EEGMMIDB S021--S040 (\(N=20\)), discrimination-based selection improved
macro AUC by 0.0561 versus covariance-NLL and 0.0920 versus Ledoit--Wolf
(both Holm \(p<0.006\)); five and four participants, respectively, suffered
losses greater than 0.02 AUC. BNCI2014\_002 (\(N=14\)) and BNCI2015\_001
(\(N=12\)) preserved positive effect directions but neither passed its
predeclared Holm family. Separately frozen Shin2017A (\(N=29\)) improved
by 0.0222 versus covariance-NLL (Holm \(p=0.0223\)), with five participants
materially harmed. A prospective representation interaction
(\(-0.0778\) AUC, Holm \(p=0.0133\)) refuted the original claim of a
universally favourable filter bank. A post-outcome whole-path sensitivity
analysis did not establish population frontier spans above 0.02 AUC.
\textbf{Conclusion:}
The regularisation \emph{selection objective} is a consequential
measurement-validity choice in the studied zero-target EEG pipelines.
The findings support neither a universal strength multiplier nor
participant-level safety or an invariably improved spectral representation.
\end{abstract}
